\documentclass[11pt]{article}
\usepackage[margin=1in]{geometry}
\usepackage{amsmath,amssymb,amsthm,mathtools}
\usepackage{booktabs}
\usepackage{microtype}
\usepackage[hidelinks]{hyperref}
\usepackage[capitalize,noabbrev]{cleveref}

\newtheorem{definition}{Definition}[section]
\newtheorem{axiom}[definition]{Axiom}
\newtheorem{assumption}[definition]{Assumption}
\newtheorem{conjecture}[definition]{Conjecture}
\newtheorem{lemma}[definition]{Lemma}
\newtheorem{proposition}[definition]{Proposition}
\newtheorem{theorem}[definition]{Theorem}
\newtheorem{corollary}[definition]{Corollary}
\theoremstyle{remark}
\newtheorem{interpretation}[definition]{Interpretation}
\newtheorem{empirical}[definition]{Empirical statement}
\newtheorem{established}[definition]{Established physical result}
\newtheorem{newresult}[definition]{New result claimed by this work}
\newtheorem{counterexample}[definition]{Counterexample}
\newtheorem{failuremode}[definition]{Failure mode}

\newcommand{\Real}{\mathbb{R}}
\newcommand{\Prob}{\mathsf{Prob}}
\newcommand{\TV}{\mathrm{TV}}
\newcommand{\id}{\mathrm{id}}

\title{Empirical Realization, Equivalence, and Limits\\
\large A Protocol-Relative Analysis for Finite-Dimensional Classical Mechanics}
\author{PRINCIPIA PHYSICA Pilot}
\date{August 2026}

\begin{document}
\maketitle

\begin{abstract}
The thesis that physics studies empirically realized compositional mathematical structures is treated here as a conjectural research program, not as an established characterization of physics.  We separate a mathematically admissible model from an empirical realization by specifying targets, preparations, interventions, measurement kernels, calibration conventions, and error models.  Observational equivalence is defined relative to such a protocol.  We prove that it is an equivalence relation and that every protocol factors through the corresponding quotient, while a bare mathematical isomorphism need not preserve empirical content.  Finite observations underdetermine smooth dynamics; limiting agreement is domain- and topology-dependent; and singular limits can invalidate unrestricted claims of reduction.  Relativistic kinetic energy and the small-frequency harmonic oscillator provide explicit counterexamples to global uniform convergence.  Finally, falsification is formulated as exclusion from a protocol-indexed acceptance region.  The results identify what must be preserved, what may be lost, and which empirical inputs are required before a mathematical construction can be called physically realized.
\end{abstract}

\section{Scope and epistemic status}

\begin{conjecture}[Founding thesis under investigation]
Physics is the study of empirically realized compositional mathematical structures.
\end{conjecture}
This sentence is not used as an axiom below.  The mathematical results test whether its central predicate, ``empirically realized,'' can be made precise without identifying mathematical equivalence with physical equivalence.

\begin{assumption}[Finite-dimensional pilot scope]\label{ass:scope}
A model has a finite-dimensional smooth state manifold $X$, a set $U$ of admissible interventions, and evolution maps $\Phi^u_{t,s}:X\to X$ whenever they are defined.  Preparations and observations are restricted to specified sets.  Field theory, infinite-dimensional dynamics, general relativity, and quantum theory are boundary tests only.
\end{assumption}
The assumption is a restriction of inquiry.  It is not a theorem about nature.  Hamiltonian models fit it when $X$ is a finite-dimensional symplectic manifold and $\Phi$ is a Hamiltonian flow; incomplete vector fields remind us that global evolution cannot be presumed.

\begin{established}[Hamiltonian mechanics]\label{est:ham}
For a smooth Hamiltonian $H$ on a symplectic manifold $(X,\omega)$, the equation
\[
 \iota_{X_H}\omega=dH
\]
defines a unique Hamiltonian vector field, and its integral curves give Hamiltonian evolution on their intervals of existence.  This is a standard result of symplectic mechanics \cite{AbrahamMarsden1978,MarsdenRatiu1999}.
\end{established}

\begin{empirical}[Status of classical mechanics]\label{emp:classical}
Finite-dimensional classical mechanics is supported within experimentally established macroscopic regimes and is known not to be universally adequate.  This paper introduces no new experimental result.  Any claim that a particular classical model realizes a target system must be supplied by target-specific calibration and data.
\end{empirical}

The distinction between models and their use in representing targets is standard in philosophy of scientific representation \cite{FriggNguyen2025}.  Van Fraassen's semantic account likewise distinguishes a family of models from empirical substructures presented as candidates for representing observable phenomena \cite{vanFraassen1980}.  Our definitions are operational refinements suited to the present mathematical scope.

\section{Mathematical models and empirical realizations}

\begin{definition}[Mathematically admissible model]\label{def:model}
A deterministic controlled model is a tuple
\[
 M=(X,U,T,\Phi,\mathcal O),
\]
where $X$ is a smooth state manifold, $U$ is a set of admissible intervention histories, $T\subseteq\Real$ is a time domain, $\Phi$ is a compatible family of evolution maps, and $\mathcal O$ is a declared class of mathematical observables $f:X\to Y_f$.  Admissibility means only that the tuple satisfies its declared mathematical typing and axioms.
\end{definition}
Thus any smooth Hamiltonian on any finite-dimensional symplectic manifold is a candidate mathematical model.  Admissibility alone supplies neither a target nor an instrument reading.

\begin{definition}[Experimental protocol]\label{def:protocol}
A protocol for a target $S$ is a tuple
\[
 P=(A,I,J,\pi,\{K_j\}_{j\in J},C,E).
\]
Here $A$ is a set of preparation labels; $I$ is a set of allowed intervention labels; $J$ is a set of measurement settings; $\pi:A\to\Prob(X)$ assigns model-state distributions to preparations; and
\[
 K_j:X\times\mathcal B(Y_j)\longrightarrow[0,1]
\]
is a Markov kernel giving predicted outcome probabilities.  The calibration record $C$ fixes units, reference standards, clock conventions, and the empirical association of labels with laboratory operations.  The error model $E$ specifies sampling error, instrument uncertainty, and a model-discrepancy allowance.
\end{definition}
For $a\in A$, intervention $i\in I$, setting $j\in J$, and time $t$, the protocol predicts
\begin{equation}\label{eq:prediction}
 p^M_{aijt}(B)=\int_X K_j\!\left(\Phi^i_{t,0}(x),B\right)\,d\pi_a(x),
 \qquad B\in\mathcal B(Y_j).
\end{equation}
A noisy instrument is represented by $K_j$ rather than being silently identified with a function on phase space.

\begin{definition}[Empirical realization]\label{def:realization}
A realization claim is a tuple $R=(S,M,P,D,\alpha)$, where $S$ is a target, $M$ an admissible model, $P$ a protocol, $D$ accepted data, and $\alpha$ an adequacy rule.  The statement ``$M$ is empirically realized by $S$ relative to $(P,D,\alpha)$'' means that (i) the associations recorded in $C$ have been operationally implemented and (ii) the distributional predictions in \eqref{eq:prediction} satisfy $\alpha(D,p^M,E)$.
\end{definition}

\begin{interpretation}[Relational and corrigible realization]\label{int:relative}
Empirical realization is a relation involving model, target, protocol, data, and tolerance.  It is not an intrinsic property of the bare tuple $M$.  New interventions, improved calibration, or a revised error model can change the warranted realization claim without changing $M$.
\end{interpretation}

\begin{counterexample}[Admissibility without realization]\label{cx:oscillator}
For every $\omega>0$, the Hamiltonian
\[
 H_\omega(q,p)=\frac{p^2}{2m}+\frac12m\omega^2q^2
\]
on $(\Real^2,dq\wedge dp)$ is mathematically admissible.  Fix a target whose measured displacement grows as $q(t)=q_0+vt$ over the tested interval with errors much smaller than $|q_0(\cos\omega t-1)+(v/\omega)(\sin\omega t-\omega t)|$.  The oscillator remains admissible but fails that realization test.  Mathematical consistency has not protected it from empirical rejection.
\end{counterexample}

\section{Observational, mathematical, and physical equivalence}

\begin{definition}[Protocol-relative observational equivalence]\label{def:obseq}
Let $M$ and $N$ be models equipped with the same empirical label sets and outcome spaces under protocol $P$.  Write $M\sim_P N$ exactly when
\[
 p^M_{aijt}=p^N_{aijt}
 \quad\text{for every admissible }(a,i,j,t).
\]
For tolerance $\varepsilon\geq0$, write $M\sim_{P,\varepsilon}N$ when
\[
 \sup_{a,i,j,t}\TV\!\left(p^M_{aijt},p^N_{aijt}\right)\leq\varepsilon.
\]
The supremum ranges only over the domain declared by $P$.
\end{definition}

\begin{proposition}\label{prop:equivrel}
Exact observational equivalence $\sim_P$ is an equivalence relation.
\end{proposition}
\begin{proof}
For every prediction index $r=(a,i,j,t)$, equality of probability measures is reflexive, symmetric, and transitive.  Applying these three properties pointwise over all $r$ proves reflexivity, symmetry, and transitivity of $\sim_P$.
\end{proof}

Approximate equivalence at fixed positive tolerance is generally not transitive: if three deterministic predictions on $\Real$ are $0$, $0.75\varepsilon$, and $1.5\varepsilon$ and distance is absolute error, the first is within $\varepsilon$ of the second and the second of the third, but the first is not within $\varepsilon$ of the third.  Consequently, approximate equivalence should be treated as a tolerance relation, not automatically as a quotient relation.

\begin{newresult}[Protocol quotient]\label{res:quotient}
Let $\mathfrak M$ be a set of models with a common protocol $P$, and let
\[
 \operatorname{Pred}_P:\mathfrak M\longrightarrow
 \prod_{a,i,j,t}\Prob(Y_j)
\]
map each model to its complete family of protocol predictions.  Then $\operatorname{Pred}_P$ factors uniquely through the quotient $q:\mathfrak M\to\mathfrak M/{\sim_P}$ as an injective map $\overline{\operatorname{Pred}}_P$.  Hence the quotient is precisely the maximal distinction among models available to $P$.
\end{newresult}
\begin{proof}
By \cref{def:obseq}, $q(M)=q(N)$ exactly when $\operatorname{Pred}_P(M)=\operatorname{Pred}_P(N)$.  Define $\overline{\operatorname{Pred}}_P([M])=\operatorname{Pred}_P(M)$.  The displayed biconditional makes this definition independent of the representative and makes the induced map injective.  Since $q$ is surjective, any factorization must assign $[M]$ the value $\operatorname{Pred}_P(M)$, proving uniqueness.
\end{proof}

This result is mathematical.  Calling the quotient ``empirical content'' is the accompanying interpretation, and its scope is limited to $P$.  Expanding $P$ can refine its equivalence classes.

\begin{definition}[Mathematical isomorphism with transported interpretation]\label{def:iso}
An isomorphism from $M$ to $N$ is a bijective smooth map $F:X_M\to X_N$ intertwining all declared mathematical structure, including
\[
 F\circ\Phi^{M,i}_{t,s}=\Phi^{N,i}_{t,s}\circ F.
\]
It transports the empirical interpretation when $F_*\pi^M_a=\pi^N_a$ and
\[
 K^M_j(x,B)=K^N_j(Fx,B)
\]
for every admitted label and measurable outcome event.
\end{definition}

\begin{proposition}[Conditional bridge from isomorphism to observation]\label{prop:bridge}
An isomorphism that transports the empirical interpretation yields $M\sim_PN$.
\end{proposition}
\begin{proof}
Insert $F_*\pi^M_a=\pi^N_a$ into the integral defining $p^N_{aijt}$.  Intertwining replaces $\Phi^{N,i}_{t,0}(Fx)$ by $F(\Phi^{M,i}_{t,0}(x))$.  Kernel compatibility then replaces $K^N_j(Fy,B)$ by $K^M_j(y,B)$.  The resulting integral is $p^M_{aijt}(B)$ for every index and event $B$.
\end{proof}

\begin{counterexample}[A bare isomorphism is insufficient]\label{cx:units}
Let two copies of $\Real^2$ carry the free-particle flow.  The symplectomorphism $F(q,p)=(\lambda q,p/\lambda)$, $\lambda>0$, is a mathematical isomorphism after the Hamiltonian parameters are transformed accordingly.  Suppose both empirical interpretations nevertheless associate the numerical coordinate $q$ with the same calibrated metre reading.  Then a state with coordinate $q=1$ predicts one metre in the first representation and $\lambda$ metres after applying $F$ in the second.  Unless calibration is transported, the two interpreted models differ observationally.  The mathematical map does not determine what laboratory operation ``one metre'' denotes.
\end{counterexample}

The literature distinguishes empirical, definitional, categorical, and interpretative notions of theoretical equivalence \cite{Weatherall2019a,Weatherall2019b,Coffey2014}.  Our result supplies a narrow preservation test: formal equivalence is physically relevant only after the interpretation and its losses are stated.

\section{Underdetermination by finite data}

\begin{newresult}[Finite sampling does not identify a smooth trajectory]\label{res:finite}
Let $t_1,\ldots,t_n$ be distinct times in an open interval $I$, and let $q:I\to\Real$ be smooth.  There are infinitely many distinct smooth trajectories $q_\lambda$ satisfying $q_\lambda(t_k)=q(t_k)$ for every $k$.
\end{newresult}
\begin{proof}
Choose a nonzero smooth function $g:I\to\Real$ and set
\[
 h(t)=g(t)\prod_{k=1}^n(t-t_k),\qquad q_\lambda(t)=q(t)+\lambda h(t).
\]
For every $k$, the product vanishes at $t_k$, so $q_\lambda(t_k)=q(t_k)$.  If $\lambda\neq\mu$, then $q_\lambda-q_\mu=(\lambda-\mu)h$ is not identically zero.  Thus the family contains infinitely many distinct smooth trajectories agreeing with all sampled values.
\end{proof}

\begin{corollary}\label{cor:finite}
Finite noiseless position samples alone cannot identify an unrestricted smooth law of motion.
\end{corollary}
\begin{proof}
Every trajectory in \cref{res:finite} yields the same finite data.  Since the trajectories differ away from the sample times, the map from smooth trajectories to those data is not injective.
\end{proof}

This is a structural non-identifiability result, not a claim that scientific inference is impossible.  Identification can be recovered relative to a restricted model class, informative interventions, derivative data, or probabilistic regularization.  Such restrictions are assumptions and must not be redescribed as consequences of the observations.

\begin{failuremode}[Auxiliary dependence]
A failed prediction may implicate dynamics, preparation, measurement kernel, clock synchronization, calibration, or error model.  A test falsifies the conjunction used to derive its prediction.  Diagnosing which component failed requires additional interventions or independent calibration.
\end{failuremode}

\begin{failuremode}[Protocol impoverishment]
If every $K_j$ is constant in $x$, all models are observationally equivalent under $P$.  The quotient theorem remains valid, but the quotient has a single class and carries no state-discriminating information.  Formal correctness of a protocol therefore does not entail empirical informativeness.
\end{failuremode}

\section{Limiting relations and their domains}

A statement ``theory $A$ is the limit of theory $B$'' is incomplete until it identifies (i) a parameter, (ii) maps between state and observable spaces, (iii) a domain, (iv) a topology or error metric, and (v) the order of limits.  Limiting agreement need not amount to identity, derivability, or global approximation; discussions of reduction and emergence emphasize this dependence on limiting structure \cite{Butterfield2011a,Butterfield2011b}.

\begin{definition}[Protocol-compatible limiting relation]\label{def:limit}
Let $M_\epsilon$ and $M_0$ share protocol labels after comparison maps have been supplied.  We say $M_\epsilon\to_P M_0$ on $D$ when
\[
 \sup_{r\in D}\TV\!\left(p^{M_\epsilon}_r,p^{M_0}_r\right)\longrightarrow0
 \quad(\epsilon\to0).
\]
Here $D$ is an explicitly declared subset of preparations, interventions, settings, and times.
\end{definition}

\begin{established}[Relativistic kinetic energy]\label{est:rel}
For a free particle of rest mass $m>0$, special relativity assigns kinetic energy
\[
 T_c(p)=\sqrt{m^2c^4+p^2c^2}-mc^2,
\]
while Newtonian mechanics assigns $T_N(p)=p^2/(2m)$.  The former has the low-momentum expansion
\[
 T_c(p)=\frac{p^2}{2m}-\frac{p^4}{8m^3c^2}+O(c^{-4})
\]
for fixed $p$ as $c\to\infty$.
\end{established}

\begin{proposition}[Compact-domain Newtonian limit]\label{prop:rel-limit}
For every $P<\infty$, $T_c\to T_N$ uniformly on $|p|\leq P$ as $c\to\infty$.  The convergence is not uniform on all of $\Real$.
\end{proposition}
\begin{proof}
Write $z=p^2/(m^2c^2)$.  Taylor's theorem for $\sqrt{1+z}$ on $0\leq z\leq P^2/(m^2c^2)$ gives a constant $C_P$ and sufficiently large $c$ such that
\[
 \sup_{|p|\leq P}|T_c(p)-T_N(p)|\leq C_Pc^{-2}.
\]
This proves compact-uniform convergence.  For nonuniformity, choose $p=mc$.  Then
\[
 |T_c(mc)-T_N(mc)|=mc^2\left|\sqrt2-1-\frac12\right|,
\]
which diverges with $c$.  Therefore the supremum over $p\in\Real$ does not tend to zero.
\end{proof}

\begin{interpretation}
The phrase ``Newtonian mechanics is the $c\to\infty$ limit'' is warranted here only for a bounded momentum regime after rest energy has been subtracted and observables have been matched.  The countersequence $p=mc$ exits that regime.
\end{interpretation}

\begin{proposition}[Finite-time but not global oscillator limit]\label{prop:osc-limit}
Let $q_\omega$ solve $\ddot q_\omega+\omega^2q_\omega=0$ with $q_\omega(0)=q_0$ and $\dot q_\omega(0)=v_0$, and let $q_0^{\mathrm{free}}(t)=q_0+v_0t$.  For every $T<\infty$,
\[
 \sup_{|t|\leq T}|q_\omega(t)-q_0^{\mathrm{free}}(t)|\to0
 \quad(\omega\to0).
\]
If $(q_0,v_0)\neq(0,0)$, convergence is not uniform on $t\in\Real$.
\end{proposition}
\begin{proof}
The solution is
\[
 q_\omega(t)=q_0\cos(\omega t)+\frac{v_0}{\omega}\sin(\omega t).
\]
For $|t|\leq T$, the bounds $|1-\cos x|\leq x^2/2$ and $|\sin x-x|\leq|x|^3/6$ give
\[
 |q_\omega(t)-q_0-v_0t|
 \leq \frac{|q_0|\omega^2T^2}{2}+\frac{|v_0|\omega^2T^3}{6},
\]
which tends to zero.  For global failure, if $v_0\neq0$, evaluate at $t=2\pi/\omega$: $q_\omega=q_0$ while $q_0^{\mathrm{free}}=q_0+2\pi v_0/\omega$.  If $v_0=0$ and $q_0\neq0$, evaluate at $t=\pi/\omega$: the discrepancy is $2|q_0|$.
\end{proof}

\begin{failuremode}[Noncommuting extrapolation]
Taking $\omega\to0$ at fixed time yields the free trajectory, whereas choosing times proportional to $1/\omega$ retains or amplifies the discrepancy.  A finite-time approximation cannot be extrapolated to an unrestricted temporal equivalence.
\end{failuremode}

\begin{failuremode}[Loss under coarse graining]
If a measurement map $m:X\to Y$ is noninjective, distinct states in a fiber $m^{-1}(y)$ are observationally indistinguishable for that measurement.  Composition can expose the lost distinction: coupling to an ancillary system may make later outcomes depend on which state in the fiber was present.  Therefore empirical equivalence under isolated measurements need not survive an enlarged intervention class.
\end{failuremode}

\section{Falsification as protocol-indexed exclusion}

\begin{definition}[Acceptance region]\label{def:accept}
For a realized-model candidate $(M,P,E)$ and test index $r$, let $\mathcal A_{M,P,E}(r)\subseteq\mathcal D_r$ be the set of data records accepted by the predeclared adequacy rule.  A datum $d_r$ falsifies the candidate relative to the test when $d_r\notin\mathcal A_{M,P,E}(r)$.
\end{definition}
This definition does not identify statistical rejection with logical proof that every component except the dynamics is correct.  It records exactly which conjunction has failed its stated test.

\begin{theorem}[Equivalence and falsification]\label{thm:falsification}
Suppose $M\sim_PN$ and both candidates use the same error model and adequacy rule, with acceptance regions determined solely by their protocol predictions.  Then
\[
 \mathcal A_{M,P,E}(r)=\mathcal A_{N,P,E}(r)
\]
for every $r$.  Consequently, no datum admitted by $P$ falsifies exactly one of $M$ and $N$.
\end{theorem}
\begin{proof}
Observational equivalence gives $p^M_r=p^N_r$ for every $r$.  The acceptance-region construction has, by hypothesis, only the prediction, common error model, and common rule as inputs.  Substitution of equal predictions therefore gives equal acceptance regions.  Membership of any datum is the same in equal sets.
\end{proof}

\begin{corollary}[Discrimination requires protocol expansion]\label{cor:expand}
To discriminate exactly observationally equivalent candidates empirically, at least one of the preparation, intervention, measurement, time, calibration, or error-model components must be expanded or revised.
\end{corollary}
\begin{proof}
Under the unchanged protocol and common assessment rule, \cref{thm:falsification} forbids differential falsification.  A discriminating test must therefore violate at least one of those hypotheses, which means changing or extending a protocol component or its assessment model.
\end{proof}

\begin{counterexample}[Equivalence does not imply truth]
Two models may assign the same probability $1/2$ to every binary outcome while repeated observations are stably concentrated near frequency $0.9$ under a validated sampling protocol.  They are observationally equivalent to each other but jointly inadequate.  Equivalence is a comparative relation, not an adequacy certificate.
\end{counterexample}

\begin{proposition}[A sufficient falsification witness]\label{prop:witness}
Let an observable prediction be a scalar $\mu_M(r)$, and suppose the adequacy rule accepts exactly data satisfying $|d_r-\mu_M(r)|\leq\delta_r$, where $\delta_r$ is fixed before observing $d_r$.  If $|d_r-\mu_M(r)|>\delta_r$, then the realization claim fails relative to that protocol and rule.
\end{proposition}
\begin{proof}
The acceptance region is the closed interval $[\mu_M(r)-\delta_r,\mu_M(r)+\delta_r]$.  The strict inequality states that $d_r$ lies outside this interval, which is falsification by \cref{def:accept}.
\end{proof}

Post hoc enlargement of $\delta_r$ can protect any finite discrepancy and therefore destroys the force of this test unless the revision is independently justified and retested.  Conversely, setting $\delta_r=0$ for a noisy instrument confuses idealized mathematics with achievable discrimination.

\section{Composition and empirical interpretation}

\begin{assumption}[Mathematical composition]\label{ass:composition}
For independently represented systems $M_1,M_2$, a composite mathematical model $M_1\otimes M_2$ is available, with declared state-space and process composition.  In elementary Hamiltonian mechanics without constraints, the state space is often $X_1\times X_2$ with symplectic form $\omega_1\oplus\omega_2$.
\end{assumption}
This assumption can fail for constrained or quotient constructions, so it is not elevated to a universal axiom.

\begin{conjecture}[Empirical compositionality]\label{conj:empcomp}
A useful compositional empirical theory should provide protocol-composition maps for which independent preparation and measurement predictions satisfy
\[
 \TV\!\left(p^{M_1\otimes M_2}_{r_1\otimes r_2},
 p^{M_1}_{r_1}\otimes p^{M_2}_{r_2}\right)\leq\varepsilon_{12}
\]
whenever the target systems are operationally isolated to the declared tolerance.  Interactions must be represented by additional morphisms rather than hidden inside the tensor symbol.
\end{conjecture}
The conjecture is testable in a specified experimental domain.  Exact factorization is not expected when preparation correlations, common-mode noise, or residual interaction exceed the error allowance.

\begin{interpretation}[What the founding thesis can mean]
Within the pilot scope, the founding thesis can be rendered as a methodological claim: physics selects, interprets, composes, and tests mathematical structures through explicit empirical protocols.  The word ``realized'' then denotes warranted adequacy relative to such protocols, not an unexplained embedding of mathematics into the world.  The formulation does not establish that all of physics is exhausted by this scheme.
\end{interpretation}

\section{Failure checklist and conclusions}

A defensible realization or equivalence claim must answer the following questions.
\begin{enumerate}
 \item What target, preparation, intervention, and measurement operations are in scope?
 \item Which maps connect laboratory records to model states and observables?
 \item What units, clocks, reference frames, and calibration standards are fixed?
 \item What stochastic and systematic errors are admitted, and were tolerances fixed before testing?
 \item Which structure is preserved by an alleged equivalence, and which distinctions are forgotten?
 \item Over which state, parameter, and time domains does a limiting estimate hold?
 \item Is convergence pointwise, compact-uniform, global-uniform, probabilistic, or merely formal?
 \item Which auxiliary assumptions participate in a failed prediction?
 \item Does an equivalence survive enlarged interventions and composition with ancillary systems?
\end{enumerate}

The mathematical conclusions are bounded.  Exact observational equivalence is an equivalence relation and determines a protocol quotient.  Isomorphism implies observational equivalence only with a compatible transported interpretation.  Finite samples do not identify unrestricted smooth trajectories.  Two familiar limits are uniform only on restricted domains.  Equivalent models share falsification status only under the same full protocol and assessment rule.  The empirical conclusion is correspondingly conditional: a mathematical structure is not physically realized until the required correspondence and adequacy claims are supported by calibration and data.  The broader claim that this framework captures the essence of physics remains conjectural.

\begin{thebibliography}{99}

\bibitem{AbrahamMarsden1978}
R. Abraham and J. E. Marsden,
\emph{Foundations of Mechanics}, 2nd ed., Benjamin/Cummings (1978); reprint, AMS Chelsea (2008).
\href{https://doi.org/10.1090/chel/364}{doi:10.1090/chel/364}.

\bibitem{Butterfield2011a}
J. Butterfield,
``Less is Different: Emergence and Reduction Reconciled,''
\emph{Foundations of Physics} \textbf{41}, 1065--1135 (2011).
\href{https://doi.org/10.1007/s10701-010-9516-1}{doi:10.1007/s10701-010-9516-1}.

\bibitem{Butterfield2011b}
J. Butterfield,
``Emergence, Reduction and Supervenience: A Varied Landscape,''
\emph{Foundations of Physics} \textbf{41}, 920--959 (2011).
\href{https://doi.org/10.1007/s10701-011-9549-0}{doi:10.1007/s10701-011-9549-0}.

\bibitem{Coffey2014}
K. Coffey,
``Theoretical Equivalence as Interpretative Equivalence,''
\emph{British Journal for the Philosophy of Science} \textbf{65}, 821--844 (2014).
\href{https://doi.org/10.1093/bjps/axt034}{doi:10.1093/bjps/axt034}.

\bibitem{FriggNguyen2025}
R. Frigg and J. Nguyen,
``Scientific Representation,''
\emph{Stanford Encyclopedia of Philosophy} (substantive revision 2025).
\url{https://plato.stanford.edu/entries/scientific-representation/}.

\bibitem{MarsdenRatiu1999}
J. E. Marsden and T. S. Ratiu,
\emph{Introduction to Mechanics and Symmetry}, 2nd ed., Springer (1999).
\href{https://doi.org/10.1007/978-0-387-21792-5}{doi:10.1007/978-0-387-21792-5}.

\bibitem{vanFraassen1980}
B. C. van Fraassen,
\emph{The Scientific Image}, Oxford University Press (1980).
\url{https://global.oup.com/academic/product/the-scientific-image-9780198244271}.

\bibitem{Weatherall2019a}
J. O. Weatherall,
``Part 1: Theoretical Equivalence in Physics,''
\emph{Philosophy Compass} \textbf{14}, e12592 (2019).
\href{https://doi.org/10.1111/phc3.12592}{doi:10.1111/phc3.12592}.

\bibitem{Weatherall2019b}
J. O. Weatherall,
``Part 2: Theoretical Equivalence in Physics,''
\emph{Philosophy Compass} \textbf{14}, e12591 (2019).
\href{https://doi.org/10.1111/phc3.12591}{doi:10.1111/phc3.12591}.

\end{thebibliography}

\end{document}
